% \documentclass{article}
% \usepackage[utf8]{inputenc}
% \usepackage{amsmath}
% \usepackage{graphicx}
% \usepackage{booktabs}
% \usepackage{array}
% \usepackage{geometry}\geometry{a4paper, margin=1in}
% \usepackage{hyperref}
% \usepackage{parskip}

% \title{Architectural Description: Human Benchmark \\ \Large Deployment View}
% \author{Generated by AI Software Architect Assistant}
% \date{\today}

% \begin{document}

% \maketitle
% \tableofcontents
% \newpage

\section{Introduction to the Deployment View}
The Deployment View of the Human Benchmark system describes the environment into
which the system will be deployed, its dependencies on other elements of that
environment.
% This view is crucial for understanding how the various software
% components of Human Benchmark are distributed across hardware nodes, how these
% nodes are connected, and what third-party software and hardware dependencies
% exist.

The Deployment View addresses the following primary concerns:
\begin{table}[H]
    \centering
    \begin{tabular}{|p{0.25\textwidth}|p{0.75\textwidth}|}
        \hline
        \textbf{Concern}         & \textbf{Description}                                                                                   \\
        \hline
        Runtime Platform         & The software and hardware environment required to run the deployment.                                  \\
        \hline
        Hardware                 & The specification of hardware or hosting required for deployment.                                      \\
        \hline
        Third-party Software     & Any third-party software dependencies for each deployment element.                                     \\
        \hline
        Technology Compatibility & Compatibility between system components and their deployment environments.                             \\
        \hline
        Network Requirements     & Network requirements and capacity for supporting multiplayer functionality.                            \\
        \hline
        Physical Constraints     & Any physical constraints that may impact deployment, such as server locations or hardware limitations. \\
        \hline
    \end{tabular}
    \caption{Primary Concerns Addressed by the Deployment View}
    \label{tab:deployment_view_concerns}
\end{table}

THe Deployment View uses these models to describe the deployment architecture:
\begin{table}[H]
    \centering
    \begin{tabular}{|p{0.25\textwidth}|p{0.75\textwidth}|}
        \hline
        \textbf{Model}              & \textbf{Description}                                                                        \\
        \hline
        Runtime Platform Model      & Describes the software and hardware environment required to run the Human Benchmark system. \\
        \hline
        Network Model               & Describes the network topology, services, and bandwidth requirements for the deployment.    \\
        \hline
        Technology Dependency Model & Lists the software dependencies required for each deployment element.                       \\
        \hline
        Intermodel relationships    & Describes how the deployment model relates to other architectural views.                    \\
        \hline
    \end{tabular}
    \caption{Models Used in the Deployment View}
    \label{tab:deployment_view_models}
\end{table}

The primary stakeholders for this view are:
\begin{table}[H]
    \centering
    \begin{tabular}{|p{0.25\textwidth}|p{0.3\textwidth}|p{0.35\textwidth}|}
        \hline
        \textbf{Stakeholder}  & \textbf{Role}                                                 & \textbf{Interest}                                                \\
        \hline
        System Administrators & Manage and maintain the deployed system                       & Ensure uptime and performance                                    \\
        \hline
        Developers            & Implement the system                                          & Understand deployment considerations                             \\
        \hline
        Testers               & Validate the system in production-like environments           & Configure appropriate test environments                          \\
        \hline
        Communicators         & Explain deployment aspects to other stakeholders              & Ensure understanding of how the system operates in production    \\
        \hline
        Assessors             & Evaluate the technical suitability of the deployment approach & Ensure it meets quality attributes like performance and security \\
        \hline
    \end{tabular}
    \caption{Stakeholders for the Deployment View of Human Benchmark}
    \label{tab:deployment_view_stakeholders}
\end{table}

% The primary stakeholders for this view are:
% \begin{itemize}
%     \item System administrators: Who will manage and maintain the deployed system
%     \item Developers: Who need to understand deployment considerations during implementation
%     \item Testers: Who need to configure appropriate test environments
%     \item Communicators: Who explain deployment aspects to other stakeholders
%     \item Assessors: Who evaluate the technical suitability of the deployment approach
% \end{itemize}

\section{Core Architectural Elements of the Deployment View}
We will be using these abreviations throughout the Deployment View:
\begin{itemize}
    \item \textbf{VM}: Virtual Machine
    \item \textbf{DB}: Database (Server)
    \item \textbf{API}: Application Programming Interface
    \item \textbf{UI}: User Interface
\end{itemize}

\subsection{Runtime Platform Model}
The runtime platform model describes the software and hardware environment
required to run the Human Benchmark system.
% , including the allocation of software elements to hardware nodes.

\subsubsection{Processing Nodes}
Before we begin, we would like to include a quick reference table to the specifications of our virtual machine
which runs on MIF OpenNebula and database which runs on Aiven PostgreSQL services.

\begin{table}[H]
    \centering
    \caption{Specifications of Virtual Machine and Database Server}
    \label{tab:vm_db_specs}
    \begin{tabular}{|p{0.3\textwidth}|p{0.2\textwidth}|p{0.2\textwidth}|p{0.2\textwidth}|}
        \hline
        \textbf{Component} & \textbf{CPU Units} & \textbf{RAM} & \textbf{Storage} \\
        \hline
        Virtual Machine    & 0.4                & 2 GB         & 20 GB            \\
        \hline
        Database Server    & 1.0                & 1 GB         & 5 GB             \\
        \hline
    \end{tabular}
\end{table}

\begin{figure}[H]
    \centering
    \includegraphics[width=0.9\textwidth]{AutomationAndSections/Sections/6deploymentView/Images/processsing.png}
    \caption{Processing Nodes and Their Relationships in the Deployment Architecture}
    \label{fig:processing_nodes_relationships}
\end{figure}

Our main processing nodes are:
\begin{itemize}
    \item \textbf{Virtual machine (VM)}: The main processing node for the Human Benchmark system, running on MIF OpenNebula.
    \item \textbf{Database Server (DB)}: Hosts the PostgreSQL database for Human Benchmark, provided by Aiven.
\end{itemize}




% TODO: DIAGRAM - Deployment architecture showing the main processing nodes and their relationships
% SUGGESTED_TYPE: UML Deployment Diagram
% KEY_ELEMENTS_TO_DEPICT: Application Servers, Web Servers, Database Servers, Client Devices, Kubernetes cluster, Load Balancers, Network connections
% STAKEHOLDER_FOCUS: System administrators and developers to understand the high-level deployment structure
% \begin{figure}[h!]
%     \centering
%     \caption{Illustrative Diagram: High-Level Deployment Architecture of Human Benchmark (Conceptual Placeholder)}
%     \label{fig:deployment_high_level_overview}
%     \textit{Comment: This diagram should illustrate the main processing nodes and their connections within the Kubernetes environment.}
% \end{figure}

\subsubsection{Runtime Containers}

Human Benchmark employs containers for simple deployment:

\begin{figure}[H]
    \centering
    \includegraphics[width=0.9\textwidth]{AutomationAndSections/Sections/6deploymentView/Images/dockerContainers.png}
    \caption{Containers Used in the Deployment of Human Benchmark}
    \label{fig:docker_containers}
\end{figure}

\begin{itemize}
    \item \textbf{Backend Container}: Contains the .NET runtime environment with
          WebSocket support, hosting the ASP.NET Core API services, user management, and multiplayer functionality. Uses Docker to execute the application.
    \item \textbf{Frontend Container}: Contains a React + Vite environment serving
          the user interface. Uses Docker to execute the application.
    \item \textbf{Database Container}: Runs on Aiven's PostgreSQL service, which is
          containerized for consistency and ease of management. We are not sure what exectution environment Aiven uses, hence the "???" in the diagram.
\end{itemize}

% The container specification includes proper resource allocation (CPU, memory) to ensure optimal performance across environments.

\subsubsection{Storage Hardware}

\begin{itemize}
    \item \textbf{Online DB Storage}: The primary data storage for saving data is
          provided by the PostgreSQL database hosted on Aiven. This storage is used
          for:
          \begin{itemize}
              \item User account data
              \item Game configurations
              \item Historical scores and performance metrics
              \item Game session data
          \end{itemize}
    \item \textbf{Backups for Online Storage}: Regular database can be enabled on Aiven. We
          have enabled daily backups with a 7-day retention policy to ensure that we
          can recover from data loss or corruption.
    \item \textbf{Online VM Storage}: The virtual machine running the Human Benchmark contains all the application containers, configuration files.
    \item \textbf{Backups for VM Storage}: We do not have a backup strategy for the VM storage. In case of a failure, we can redeploy the containers from the source code repository.
\end{itemize}

\subsubsection{Network Links}

The Human Benchmark system requires the following network connections:

\begin{itemize}
    \item \textbf{Client-to-VM JSON}: High-bandwidth, low-latency connections for responsive user interface updates.
    \item \textbf{Client-to-VM WebSocket}: Optimized for real-time multiplayer functionality, requiring persistent connections with minimal latency.
    \item \textbf{VM-to-DB SQL}: Secure connections to the PostgreSQL database for data access and storage.
\end{itemize}

\subsubsection{Runtime Element-to-Node Mapping}

The mapping of functional elements to processing nodes is as follows:

% TODO: DIAGRAM - Element-to-Node mapping showing which functional elements are deployed on which nodes
% SUGGESTED_TYPE: UML Deployment Diagram with artifacts
% KEY_ELEMENTS_TO_DEPICT: Containers (Frontend, Backend, Database), Functional elements (UI, Game Logic, User Management, Multiplayer, Data Access), Mapping relationships
% STAKEHOLDER_FOCUS: Developers and system administrators to understand where each component runs
\begin{figure}[H]
    \centering
    \includegraphics[width=0.9\textwidth]{AutomationAndSections/Sections/6deploymentView/Images/servicetocontainer.png}
    \caption{Service-to-Container Mapping in the Deployment Architecture}
    \label{fig:service_to_container_mapping}
\end{figure}

\begin{itemize}
    \item \textbf{User Interface Components} $\rightarrow$ Frontend Container (VM)
    \item \textbf{Game Logic Services} $\rightarrow$ Backend Container (VM)
    \item \textbf{User Management Services} $\rightarrow$ Backend Container (VM)
    \item \textbf{Multiplayer System Services} $\rightarrow$ Backend Container (VM) with WebSocket optimization
    \item \textbf{Data Access Services} $\rightarrow$ Backend Container (VM)
    \item \textbf{Database} $\rightarrow$ Database Container (DB)
\end{itemize}

\subsubsection{Hardware and Hosting Requirements}

For optimal performance and scalability, the VM, DB and client devices must
specific hardware and hosting requirements. From previous work (lab 2), we found
the minimal and recommended hardware requirements for VM and DB using various tools, such as
the \texttt{htop} command, \texttt{docker stats} command, and we also tested the
database performance using database load testing tools. For client devices, we
will just use the minimum requirements for modern web browsers.
Below are the tables with the hardware requirements:


\textbf{Virtual Machine (VM)}:

\begin{table}[H]
    \centering
    \begin{tabular}{|c|c|c|}
        \hline
        \textbf{Resource Type} & {Minimal needs} & {Recommended} \\
        \hline
        Storage                & 6 GB            & 10 GB         \\
        \hline
        RAM                    & 1 GB            & 2 GB          \\
        \hline
        CPU                    & 0.4 core        & 1 core        \\
        \hline
        Network Upload         & 629 Mbit/s      & 1 Gbit/s      \\
        \hline
        Network Download       & 8 Mbit/s        & 200 Mbit/s    \\
        \hline
        Network Latency        & ≈30 ms          & ≈10 ms        \\
        \hline
    \end{tabular}
    \caption{VM Resource Requirements}
    \label{tab:resource-types}
\end{table}


\textbf{Database (DB)}:

\begin{table}[H]
    \centering
    \begin{tabular}{|c|c|c|}
        \hline
        \textbf{Requirement Type} & \textbf{Minimal Requirements} & \textbf{Recommended Requirements} \\
        \hline
        Database CPU              & 1 core                        & 1 core                            \\
        \hline
        Database RAM              & 1 GB                          & 1 GB                              \\
        \hline
        Database Storage          & 1 GB                          & 5 GB                              \\
        \hline
    \end{tabular}
    \caption{DB Resource Requirements}
    \label{tab:system-requirements}
\end{table}


\textbf{Client Devices}:

\begin{table}[H]
    \centering
    \begin{tabular}{|p{0.15\textwidth}|p{0.65\textwidth}|}
        \hline
        \textbf{Device Type} & \textbf{Minimum Requirements}                      \\
        \hline
        Desktop/Laptop       & Modern web browser (Chrome, Firefox, Safari, Edge) \\
                             & 2 GHz dual-core processor                          \\
                             & 4 GB RAM                                           \\
                             & 1280 x 720 screen resolution                       \\
        \hline
        Mobile Devices       & Modern mobile browser or WebView                   \\
                             & iOS 14+ or Android 8+                              \\
                             & 1.5 GB RAM                                         \\
                             & 320dp minimum width                                \\
        \hline
    \end{tabular}
    \caption{Client Device Requirements}
    \label{tab:client_requirements}
\end{table}

\subsection{Network Model}

\subsubsection{Network Topology}

The Human Benchmark system requires a network architecture that supports reliable communication between servers and clients with minimal latency, particularly for real-time multiplayer gaming features.

% TODO: DIAGRAM - Network topology showing interconnections between system components
% SUGGESTED_TYPE: Network Diagram
% KEY_ELEMENTS_TO_DEPICT: External client connections, Load balancers, Internal network between services, Firewalls, WebSocket connections, Database connections
% STAKEHOLDER_FOCUS: Network administrators and system architects to understand connectivity requirements
\begin{figure}[H]
    \centering
    \includegraphics[width=0.7\textwidth]{AutomationAndSections/Sections/6deploymentView/Images/networkmodel.png}
    \caption{Network Topology of the Human Benchmark System}
    \label{fig:network_topology}
\end{figure}

\subsubsection{Network Services}

The following network services are required:

\begin{itemize}
    % \item \textbf{Load Balancing}: Distributes incoming traffic across multiple instances of the frontend and backend containers.
    \item \textbf{Firewalls}: Protects the system from unauthorized access and provides a security boundary.
          % \item \textbf{CDN (Content Delivery Network)}: Optional but recommended for static content delivery to improve performance for geographically distributed users.
          % \item \textbf{WebSocket Proxies}: Optimized for handling WebSocket connections required for multiplayer functionality.
\end{itemize}

\subsubsection{Bandwidth and Latency Requirements}

\begin{table}[H]
    \centering
    \caption{Network Bandwidth and Latency Requirements}
    \label{tab:network_requirements}
    \begin{tabular}{|p{0.3\textwidth}|p{0.3\textwidth}|p{0.35\textwidth}|}
        \hline
        \textbf{Connection Type} & \textbf{Bandwidth Requirement} & \textbf{Maximum Acceptable Latency} \\
        \hline
        Client-to-VM             & 1-5 Mbps per user              & 100ms (Regular interface)           \\
                                 &                                & 50ms (Reaction time tests)          \\
        \hline
        VM-to-DB                 & 700 Mbps  network              & 25ms                                \\
        \hline
    \end{tabular}
\end{table}

\subsection{Technology Dependency Model}

\subsubsection{Backend Container Dependencies}

\begin{table}[H]
    \centering
    \footnotesize % or \footnotesize if \scriptsize is too small
    \begin{tabularx}{\textwidth}{|X|c|p{0.3\textwidth}|}
        \hline
        \textbf{Technology}                     & \textbf{Version} & \textbf{Notes}                                \\
        \hline
        .NET SDK                                & 9.0.100          & API, Entity Framework Core, SignalR, Identity \\
        \hline
        ASP.NET Core Runtime                    & 9.0              & Required for hosting ASP.NET services         \\
        \hline
        Microsoft.EntityFrameworkCore           & 9.0.0            & ORM for data access                           \\
        \hline
        Npgsql.EntityFrameworkCore.PostgreSQL   & 9.0.1            & PostgreSQL EF Core provider                   \\
        \hline
        Microsoft.EntityFrameworkCore.InMemory  & 9.0.0            & For in-memory testing                         \\
        \hline
        Microsoft.EntityFrameworkCore.SqlServer & 9.0.0            & Optional SQL Server provider                  \\
        \hline
        Fleck                                   & 1.2.0            & WebSocket server support                      \\
        \hline
        Newtonsoft.Json                         & 13.0.3           & JSON serialization support                    \\
        \hline
        Microsoft.AspNetCore.Mvc.NewtonsoftJson & 8.0.8            & ASP.NET integration with Newtonsoft           \\
        \hline
        Swashbuckle.AspNetCore                  & 6.4.0            & Swagger API documentation                     \\
        \hline
        uldahlalex.websocket.boilerplate        & 2.3.4            & WebSocket boilerplate                         \\
        \hline
        Microsoft.VisualStudio.Azure.
        Containers.Tools.Targets                & 1.19.5           & Docker build tools                            \\
        \hline
        Docker                                  & 20.10.24+dfsg1   & Containerization platform                     \\
        \hline
        NUnit                                   & 3.13.3           & Unit testing framework                        \\
        \hline
        NUnit3TestAdapter                       & 4.2.1            & NUnit test integration                        \\
        \hline
        coverlet.collector                      & 6.0.0            & Code coverage tool                            \\
        \hline
        Moq                                     & 4.20.72          & Mocking framework for unit tests              \\
        \hline
        Microsoft.NET.Test.Sdk                  & 17.6.0           & Test project runner                           \\
        \hline
    \end{tabularx}
    \caption{Backend Container Software Dependencies}
    \label{tab:backend_dependencies}
\end{table}


\subsubsection{Frontend Container Dependencies}

\begin{table}[H]
    \centering
    \begin{tabular}{|c|c|c|}
        \hline
        \textbf{Technology} & \textbf{Version} & \textbf{Notes}                           \\
        \hline
        Node.js             & v22.9.0          & Runtime environment for frontend tooling \\
        \hline
        npm                 & Bundled          & Comes with Node.js                       \\
        \hline
        Vite                & 5.4.18           & Frontend build tool                      \\
        \hline
        React               & 18.3.1           & Frontend framework                       \\
        \hline
        React DOM           & 18.3.1           & DOM rendering for React                  \\
        \hline
        TypeScript          & 5.5.3            & Static typing                            \\
        \hline
        React Router DOM    & 6.26.2           & Client-side routing                      \\
        \hline
        Redux               & N/A              & Assumed used but not listed directly     \\
        \hline
        Axios               & 1.7.7            & HTTP client                              \\
        \hline
        react-use-websocket & 4.13.0           & WebSocket client integration             \\
        \hline
        ESLint              & 9.9.0            & Linting and code quality                 \\
        \hline
        Vitest              & 3.0.9            & Unit testing framework                   \\
        \hline
        nginx               & 1.24.x+          & For static hosting and reverse proxy     \\
        \hline
    \end{tabular}
    \caption{Frontend Container Software Dependencies}
    \label{tab:frontend_dependencies}
\end{table}

\subsubsection{Database Dependencies}

\begin{table}[H]
    \centering
    \begin{tabular}{|c|c|c|}
        \hline
        \textbf{Component} & \textbf{Version} & \textbf{Notes}                    \\
        \hline
        PostgreSQL         & 14.x+            & Main database engine              \\
        % \hline
        % PostGIS Extension & Optional & For location-based queries \\
        % \hline
        % pgcrypto Extension & Optional & For encryption and UUIDs \\
        \hline
        Aiven Tools        & Provided         & Backup and point-in-time recovery \\
        \hline
    \end{tabular}
    \caption{Database Software Dependencies}
    \label{tab:database_dependencies}
\end{table}

\subsubsection{CI/CD Infrastructure Dependencies}

\begin{table}[H]
    \centering
    \begin{tabular}{|p{0.3\textwidth}|p{0.7\textwidth}|}
        \hline
        \textbf{Component} & \textbf{Details}                                                     \\
        \hline
        Source Control     & GitLab repository                                                    \\
        \hline
        CI/CD System       & GitLab CI/CD with custom runner                                      \\
        \hline
        Runner Details     & Named \texttt{VMmif}, shell executor, Linux OS, version 17.9.1       \\
        \hline
        Container Engine   & Docker 20.10.24+dfsg1 (version can be higher)                        \\
        \hline
        VM Platform        & OpenNebula 4.14.2 (used to manage VM: CPU, RAM, storage, networking) \\
        \hline
        Deployment Method  & Docker-based deployment via GitLab pipeline on hosted VM             \\
        \hline
    \end{tabular}
    \caption{CI/CD Infrastructure Dependencies}
    \label{tab:cicd_dependencies}
\end{table}

\subsection{Intermodel Relationships}
\subsubsection{Relationships Between Deployment View Models}

The three models of the Deployment View—runtime platform model, network model, and technology dependency model—form a cohesive yet distinct set of related perspectives that serve different stakeholders at different project phases.

% \begin{figure}[H]
%     \centering
%     \includegraphics[width=0.7\textwidth]{AutomationAndSections/Sections/6deploymentView/Images/intermodel_relationship.png}
%     \caption{Layered Relationship Between Deployment View Models}
%     \label{fig:deployment_view_models_relationship}
% \end{figure}

The runtime platform model serves as the core of our Deployment View, providing a high-level perspective on processing nodes, containers, storage, and element mappings. It is primarily used early in the project by deployment planners and architects to make fundamental decisions about hosting environments and infrastructure requirements.

The network model can be viewed as a lower layer supporting the runtime platform, detailing the network infrastructure required to enable communication between the components defined in the runtime platform model. Network specialists use this model to ensure appropriate bandwidth, latency, and connectivity between Human Benchmark components, particularly the WebSocket connections critical for multiplayer functionality.

The technology dependency model represents a more detailed layer on top of the runtime platform model, specifying the exact software dependencies and versions required on each node. System administrators refer to this model during installation planning and configuration management closer to deployment time, ensuring that all required libraries, frameworks, and tools are properly installed and configured.

While these models are presented separately to serve different stakeholders, they maintain strong interrelationships:

\begin{itemize}
    \item Changes to the runtime platform model (e.g., adding new container types) typically require corresponding updates to network requirements and technology dependencies
    \item Network capacity decisions directly impact the performance capabilities of the runtime platform
    \item Technology choices constrain both the runtime platform options and network requirements
\end{itemize}

For Human Benchmark, these relationships are particularly important when considering scaling strategies, as increased user load affects all three models simultaneously—requiring additional runtime resources, network capacity, and potentially new technology dependencies to maintain performance.

\section{Perspective Priorities for the Deployment View}

\begin{table}[H]
    \centering
    \caption{Perspective Priorities for the Deployment View of Human Benchmark}
    \label{tab:perspective_priorities_deployment}
    \begin{tabular}{p{0.3\textwidth}p{0.15\textwidth}p{0.45\textwidth}}
        % \toprule
        \textbf{Architectural Perspective} & \textbf{Priority} & \textbf{Rationale for Human Benchmark within this View}                                                                                                                              \\
        % \midrule
        Accessibility                      & \greencircle  Low               & While accessibility is important for the UI, it has minimal impact on deployment architecture decisions.                                                                             \\
        % \midrule
        Availability and Resilience        & \redcircle  High              & Critical for a system with multiplayer features and competitive scoring. Deployment must ensure high availability through redundancy and proper failover mechanisms.                 \\
        % \midrule
        Development Resource               & \yellowcircle  Medium            & The containerized architecture simplifies development resource allocation. \\
        % \midrule
        Evolution                          & \yellowcircle  Medium            & The container-based deployment facilitates future evolution, allowing components to be updated independently as the system grows.                                                    \\
        % \midrule
        Internationalization               & \greencircle  Low               & Has minimal impact on deployment architecture, primarily affecting content and UI layers.                                                                                            \\
        % \midrule
        Location                           & \yellowcircle  Medium            & Geographic distribution of users impacts latency for reaction time tests; CDN and strategically located server deployments can mitigate this.                                        \\
        % \midrule
        Performance and Scalability        & \redcircle  High              & Critical for a reaction time testing system where millisecond precision matters. The deployment must support scaling to accommodate varying loads and ensure consistent performance. \\
        % \midrule
        Regulation                         & \greencircle  Low               & Human Benchmark primarily deals with non-sensitive user data, though basic compliance with data protection regulations affects storage location decisions.                           \\
        % \midrule
        Security                           & \redcircle  High              & Protecting user accounts and ensuring fair gameplay requires secure deployment architecture with proper network segmentation and access controls.                                    \\
        % \midrule
        Usability                          & \greencircle  Low               & Primarily affects UI design rather than deployment architecture.                                                                                                                     \\
        % \bottomrule
    \end{tabular}
\end{table}

\section{Impact of Key Perspectives on the Deployment View}

\subsection{Availability and Resilience Perspective}

Applying the Availability and Resilience perspective to the Deployment View of Human Benchmark significantly shapes its architecture in the following ways:

\begin{itemize}
    \item \textbf{Redundant Components}: The deployment includes multiple instances of frontend and backend containers to prevent single points of failure. This is reflected in the runtime platform model's specification of minimum node counts and high availability configurations.

    \item \textbf{Database Replication}: The PostgreSQL database hosted on Aiven is configured with standby replicas to ensure data availability in case of primary database failure. This is captured in the database hosting specifications.

    \item \textbf{Kubernetes Orchestration}: The use of Kubernetes for container orchestration enables automatic recovery from node or container failures, enhancing system resilience. This is a core aspect of the runtime platform model.

    \item \textbf{Backup Strategy}: Regular automated backups with retention policies protect against data loss, as specified in the database hosting section.

    \item \textbf{Network Redundancy}: The network model includes redundant paths and load balancing to prevent network-related outages, particularly important for the WebSocket connections that enable multiplayer functionality.
\end{itemize}

\subsection{Performance and Scalability Perspective}

The Performance and Scalability perspective profoundly influences the Deployment View in these ways:

\begin{itemize}
    \item \textbf{Horizontal Scaling}: The runtime platform model specifies a minimum number of worker nodes with the ability to scale horizontally based on load, ensuring the system can handle varying user counts without performance degradation.

    \item \textbf{Resource Allocation}: Detailed specifications for CPU, memory, and storage in the hardware requirements section ensure adequate resources for performance-critical operations, particularly for reaction time tests where millisecond precision is needed.

    \item \textbf{Network Optimization}: The network model specifies strict latency requirements, especially for WebSocket connections used in multiplayer games and reaction time tests. This directly addresses the system's core functionality requirements.

    \item \textbf{Content Delivery Network}: The optional CDN recommendation helps distribute static content closer to users, reducing load times and improving the user experience globally.

    \item \textbf{Database Performance}: The database hosting specifications include performance-oriented configurations to ensure rapid data access, essential for recording and retrieving high-precision timing data.
\end{itemize}

\subsection{Security Perspective}

Applying the Security perspective to the Deployment View results in these architectural decisions:

\begin{itemize}
    \item \textbf{Network Segmentation}: The network model includes firewalls and defines security boundaries between public-facing and internal components, protecting sensitive parts of the system.

    \item \textbf{Authentication Components}: The technology dependency model identifies the authentication libraries and mechanisms used to secure user access.

    \item \textbf{Database Security}: The database configuration includes security extensions (pgcrypto) and proper access controls to protect user data and prevent score manipulation.

    \item \textbf{Secure Communications}: The network model specifies secure communication protocols for all connections, particularly for user authentication and multiplayer session management.

    \item \textbf{Container Security}: The deployment approach includes container security considerations, ensuring that containerized components operate with minimum required privileges.
\end{itemize}

\section{Importance and Justification of the Deployment View for Human Benchmark}

The Deployment View is of critical importance for the Human Benchmark system due to several factors that directly affect its core functionality and quality attributes:

\subsection{Precision Timing Requirements}

Human Benchmark's primary purpose is to measure and test users' reaction times and cognitive abilities with millisecond precision. The deployment architecture directly impacts the system's ability to deliver accurate measurements:

\begin{itemize}
    \item Network latency must be minimized to ensure accurate timing measurements
    \item Processing resources must be sufficient to handle timing calculations without introducing delays
    \item Server placement and content delivery optimization are essential for global user fairness
\end{itemize}

Without a well-designed deployment architecture, the system could produce inconsistent or inaccurate reaction time measurements, undermining its core functionality.

\subsection{Multiplayer Experience Quality}

The real-time multiplayer features require specific deployment considerations:

\begin{itemize}
    \item WebSocket connections must be properly supported and optimized
    \item Network topology must facilitate low-latency communication between players
    \item Session management must be resilient to maintain fair competition
\end{itemize}

The Deployment View ensures these requirements are addressed through appropriate technology choices and resource allocations.

\subsection{Scaling for Variable User Loads}

Human Benchmark experiences variable user loads that demand a flexible deployment approach:

\begin{itemize}
    \item Containerization and Kubernetes orchestration provide scaling capabilities
    \item Resource allocation strategies must balance cost efficiency with performance
    \item Database access patterns require optimization for concurrent user sessions
\end{itemize}

The Deployment View's specification of container orchestration and scaling strategies is essential for maintaining consistent performance under varying loads.

\subsection{Stakeholder Value}

This view provides critical information to key stakeholders:

\begin{itemize}
    \item System administrators gain a clear understanding of what needs to be deployed and maintained
    \item Developers understand the runtime environment constraints their code must operate within
    \item Testers can configure appropriate test environments that accurately represent production
    \item Assessors can evaluate whether the deployment approach will deliver required quality attributes
\end{itemize}

Without this view, these stakeholders would lack essential information for their roles, potentially leading to deployment failures, performance issues, or incompatibilities.

\section{Cross-View Consistency Considerations}

The Deployment View must maintain consistency with other architectural views of the Human Benchmark system:

\subsection{Consistency with Functional View}
The runtime element-to-node mapping section of this Deployment View must align with the functional elements defined in the Functional View. Each functional element from that view must have a clear deployment location specified in this view.

\subsection{Consistency with Information View}
The database deployment specifications must support the data structures and access patterns defined in the Information View. The storage capacity, performance characteristics, and replication strategies should align with the data volume and criticality described there.

\subsection{Consistency with Concurrency View}
If a Concurrency View exists, the deployment environment must provide sufficient resources to support the concurrent processes and threads described in that view, particularly for handling multiple simultaneous users in multiplayer scenarios.

\subsection{Consistency with Development View}
The technology dependency model must align with the development environment and frameworks specified in the Development View. Any production deployment constraints that affect development decisions (such as supported libraries or technologies) must be consistent across both views.

\subsection{Maintaining Consistency}
To ensure ongoing consistency between views as the architecture evolves:

\begin{itemize}
    \item Changes to functional elements must be reflected in their deployment mappings
    \item Scaling decisions in deployment must consider functional and concurrency impacts
    \item Technology upgrades must be coordinated across development and deployment views
    \item Testing should validate cross-view consistency before deployment
\end{itemize}

This Deployment View is one part of a comprehensive Architectural Description of the Human Benchmark system. It should be reviewed and updated whenever significant changes are made to other views to maintain a coherent overall architecture.

% \end{document}